\statement{B.77}{Proved}{The four transported triple relations do not impose uniform kernel normality}
The four actual triple-point relations involving the missing retained
meridian in the ratio pencil, together with the relation at the chosen
special fiber, do not by themselves make the special-fiber deletion
kernel normal under that missing meridian. This failure holds for the
actual subgroup images in both the discrete partial presentation and
its full profinite completion. It is a statement about the specified
partial presentation; the finite certificate below does not extend
across the actual retained base monodromy.

\begin{proof}
Use the common based marking of
\path{work_DE/next_marginal_source/pencil_triple_transports.json}.
Write the generic fiber group in the special basis as
\[
 F_{20}=H*\langle b\rangle,\qquad
 H=F(x_i:1\leq i\leq20,\ i\ne11),\qquad
 b=x_{11},\quad c=x_{10}.
 \tag{B.77.1}
\]
Thus the original generic generator with index ten is \(cb^{-1}\).
The deleted peripheral indices are
\(\mathcal U=\{1,2,6,9,12,15,19,20\}\). Put
\(N=\langle\!\langle x_i:i\in\mathcal U\rangle\!\rangle_H\).
The quotient \(H/N\) is free of rank eleven. In the following table,
a string such as \(789\) denotes the literal word \(x_7x_8x_9\):
\[
\begin{array}{c|c|c}
 i&P_i&Q_i\\\hline
 1&(789)^{-1}x_6(789)&x_{10}x_{12}x_{10}^{-1}\\
 2&(2345789)^{-1}x_1(2345789)&
       (10\,12\,14\,15\,16\,17\,18)x_{19}
       (10\,12\,14\,15\,16\,17\,18)^{-1}\\
 3&x_{16}x_{20}x_{16}^{-1}&x_9^{-1}x_2x_9\\
 4&x_{15}&x_9
\end{array}
 \tag{B.77.2}
\]
These eight words and the eleven retained generators form a Nielsen
basis of \(H\). Indeed, the original deleted generators are solved
successively in the order
\(9,15,20,2,6,12,1,19\); at each step the conjugating word uses
only retained generators and deleted generators already solved.
The exact substitutions and their two inverse compositions are
included in the certificate. In particular,
\[
 A=\langle c,P_1,Q_1,\ldots,P_4,Q_4\rangle\simeq F_9,
 \qquad H=A*F_{10},
 \tag{B.77.3}
\]
where the retained generator \(x_3\) belongs to the displayed
complementary factor.

\emph{The exact partial group.}\par
Conjugate each actual local triple by its recorded transporter in
\(H\), and order its meridians by their actual local positions.
The ordered products become \(P_iQ_i b\) for \(i=1,2,3\),
and \(bP_4Q_4\) for \(i=4\). For precision, the two remaining
local meridians, besides \(b\), are respectively
\[
\begin{array}{c|c|c}
1&b^{-1}P_1b&b^{-1}Q_1b\\
2&b^{-1}Q_2^{-1}P_2Q_2b&b^{-1}Q_2b\\
3&b^{-1}Q_3b&P_3\\
4&Q_4&bP_4b^{-1}.
\end{array}
 \tag{B.77.4}
\]
A local triple relation makes its ordered product commute with
all three local meridians. Put \(R_i=P_iQ_i\). Commutation with
\(b\) first gives \([R_i,b]=1\). Using (B.77.4), the other two
relations are then exactly
\[
 bP_ib^{-1}=R_i^{-1}P_iR_i,
 \qquad bQ_ib^{-1}=R_i^{-1}Q_iR_i.
 \tag{B.77.5}
\]
Conversely (B.77.5) implies all the original local relations,
including \([R_i,b]=1\). Together with the retained-node
relation \([c,b]=1\), these are the only relations imposed in
\(G_{\mathrm{part}}\), by definition.

The map \(\eta:A\to A\) fixing \(c\) and conjugating each
rank-two free factor \(\langle P_i,Q_i\rangle\) by
\(R_i^{-1}\) is an automorphism. Consequently the exact partial
presentation is
\[
 G_{\mathrm{part}}
 =\langle H,b\mid bab^{-1}=\eta(a),\ a\in A\rangle
 \simeq(A\rtimes_{\eta}\mathbb Z)*F_{10}.
 \tag{B.77.6}
\]
For example, \(k=x_3x_9x_3^{-1}\) lies in \(N\), while
\(bkb^{-1}\) lies outside \(H\) by the reduced free-product
normal form in (B.77.6). Thus normalization of each designated
rank-two peripheral subgroup does not imply normalization of
its entire normal closure under \(H\).

\emph{A finite separating certificate.}\par
There is also an explicit finite verification requiring no
completion theorem for the HNN presentation. In the special basis set
\[
 b\longmapsto(34),\qquad x_9\longmapsto(12),\qquad
 x_3\longmapsto(13),\qquad x_i\longmapsto1
 \quad(i\notin\{3,9,11\}).
 \tag{B.77.7}
\]
Every relation in (B.77.4)--(B.77.5) and \([c,b]=1\) holds.
The image of \(H\) is the symmetric group on
\(\{1,2,3\}\), whereas the full image is \(S_4\). But
\[
             k\longmapsto(23),\qquad
             bkb^{-1}\longmapsto(24).
 \tag{B.77.8}
\]
The latter is outside the source image. The continuous extension
of this finite quotient separates \(bkb^{-1}\) from the closure
of \(H\), hence also from the closure of its deletion-kernel
image, in \(\widehat{G}_{\mathrm{part}}\). This proves the
stated full profinite version using actual images.

The exact checker is
\path{work_CE/code/check_four_triple_normality_gap.py}; its genuine
output and complete certificate are
\path{work_CE/code/check_four_triple_normality_gap.out} and
\path{work_CE/four_triple_normality_gap.json}.
It verifies both Nielsen substitutions, every literal transported
local word, all partial relations, and the orders \(6\) and
\(24\), and records hashes of the input data.

Finally the checker applies the two actual retained base
monodromies to (B.77.7), after converting back to the original
generic basis. The monodromy named \texttt{base2} fixes the tuple.
The monodromy named \texttt{base1} changes entries \(3,10,11\),
and no element of \(S_4\) conjugates the original tuple to the
resulting tuple. Among the other twenty-five filled lassos, exactly the two named
\texttt{discriminant\_2\_plus} and
\texttt{discriminant\_2\_minus} also fail; the other twenty-three
fix the tuple. These two half twists identify the two branches of
one retained hyperbola. Thus (B.77.7) is not a representation of the
actual total pencil. The remaining filled relations and retained
base actions cannot be discarded in an argument for the actual
source image. No failure of uniform normality in that full group,
and no nonidentity element of \(\mathcal K\), follows here.
\end{proof}
